\subsection*{Problema 04}

\textbf{Enunciado:} \\
Calcular la integral indefinida:
\begin{align*}
\int x^3 (1 + x^2)^{-\frac{1}{2}} \, dx
\end{align*}

\textbf{Solución:} \\
Siguiendo las indicaciones, realizamos
la sustitución sugerida. Definimos:
\begin{align*}
u = \sqrt{1 + x^2}
\end{align*}

Elevando al cuadrado ambos lados y
despejando $x^2$, obtenemos:
\begin{align*}
u^2 = 1 + x^2 \implies x^2 = u^2 - 1
\end{align*}

Ahora, diferenciamos ambos lados de la
expresión para encontrar $dx$:
\begin{align*}
2u \, du = 2x \, dx \implies x \, dx = u \, du
\end{align*}

Reescribimos la integral original
separando un factor $x$ para usar $dx$:
\begin{align*}
I &= \int x^2 (1 + x^2)^{-\frac{1}{2}} (x \, dx)
\end{align*}

Sustituimos $x^2 = u^2 - 1$,
$1 + x^2 = u^2$ y $x \, dx = u \, du$:
\begin{align*}
I &= \int (u^2 - 1) (u^2)^{-\frac{1}{2}} (u \, du)
\end{align*}

Simplificamos el término algebraico:
\begin{align*}
(u^2)^{-\frac{1}{2}} = \dfrac{1}{\sqrt{u^2}} = \dfrac{1}{u}
\end{align*}

Sustituimos esto en la integral:
\begin{align*}
I &= \int (u^2 - 1) \left(\dfrac{1}{u}\right) u \, du \\
I &= \int (u^2 - 1) \, du
\end{align*}

Integramos término a término respecto a $u$:
\begin{align*}
I &= \int u^2 \, du - \int 1 \, du \\
I &= \dfrac{u^3}{3} - u + C
\end{align*}

Finalmente, regresamos a la variable
original $x$ sustituyendo $u = \sqrt{1 + x^2}$:
\begin{align*}
I &= \dfrac{(\sqrt{1 + x^2})^3}{3} \\
  &\quad - \sqrt{1 + x^2} + C
\end{align*}

Expresando el resultado final de
forma compacta:
\begin{align*}
I &= \dfrac{(1 + x^2)^{\frac{3}{2}}}{3} \\
  &\quad - (1 + x^2)^{\frac{1}{2}} + C
\end{align*}

\begin{figure}[H]
	\centering
	\includegraphics[width=\columnwidth]{semana_05/graficas/SEM05_P04_grafica.png}
	\caption{Representación gráfica de $f(x)$ y su antiderivada $F(x)$ evaluada para la constante de integración $C = 0$ (Problema 04).}
	\label{fig:problema04}
\end{figure}
